\documentclass[a4paper]{report}


\title{Analysis}

\def\ncourse {Analysis I}

\RequirePackage{etex}
\makeatletter
\ifx \nauthor\undefined
\else
\fi

\author{Enuma Alish}
\date{\today}

\usepackage{alltt}
\usepackage{amsfonts}
\usepackage{amsmath}
\usepackage{amssymb}
\usepackage{amsthm}
\usepackage{booktabs}
\usepackage{caption}
\usepackage{enumitem}
\usepackage{fancyhdr}
\usepackage{graphicx}
\usepackage{mathdots}
\usepackage{mathtools}
\usepackage{microtype}
\usepackage{multirow}
\usepackage{pdflscape}
\usepackage{pgfplots}
\usepackage{siunitx}
\usepackage{textcomp}
\usepackage{slashed}
\usepackage{tabularx}
\usepackage{tikz}
\usepackage{tkz-euclide}
\usepackage[normalem]{ulem}
\usepackage[all]{xy}
\usepackage{imakeidx}

\makeindex[intoc, title=Index]
\indexsetup{othercode={\lhead{\emph{Index}}}}


\pgfplotsset{compat=1.12}

\ifx \ncoursehead \undefined
\def\ncoursehead{\ncourse}
\fi


\usetikzlibrary{arrows.meta}
\usetikzlibrary{decorations.markings}
\usetikzlibrary{decorations.pathmorphing}
\usetikzlibrary{positioning}
\usetikzlibrary{fadings}
\usetikzlibrary{intersections}
\usetikzlibrary{cd}

\newcommand*{\Cdot}{{\raisebox{-0.25ex}{\scalebox{1.5}{$\cdot$}}}}
\newcommand {\pd}[2][ ]{
  \ifx #1 { }
    \frac{\partial}{\partial #2}
  \else
    \frac{\partial^{#1}}{\partial #2^{#1}}
  \fi
}
\ifx \nhtml \undefined
\else
  \renewcommand\printindex{}
  \DisableLigatures[f]{family = *}
  \let\Contentsline\contentsline
  \renewcommand\contentsline[3]{\Contentsline{#1}{#2}{}}
  \renewcommand{\@dotsep}{10000}
  \newlength\currentparindent
  \setlength\currentparindent\parindent

  \newcommand\@minipagerestore{\setlength{\parindent}{\currentparindent}}
  \usepackage[active,tightpage,pdftex]{preview}
  \renewcommand{\PreviewBorder}{0.1cm}

  \newenvironment{stretchpage}%
  {\begin{preview}\begin{minipage}{\hsize}}%
    {\end{minipage}\end{preview}}
  \AtBeginDocument{\begin{stretchpage}}
  \AtEndDocument{\end{stretchpage}}

  \newcommand{\@@newpage}{\end{stretchpage}\begin{stretchpage}}

  \let\@real@section\section
  \renewcommand{\section}{\@@newpage\@real@section}
  \let\@real@subsection\subsection
  \renewcommand{\subsection}{\@ifstar{\@real@subsection*}{\@@newpage\@real@subsection}}
\fi
\ifx \ntrim \undefined
\else
  \usepackage{geometry}
  \geometry{
    papersize={379pt, 699pt},
    textwidth=345pt,
    textheight=596pt,
    left=17pt,
    top=54pt,
    right=17pt
  }
\fi

\ifx \nisofficial \undefined
\let\@real@maketitle\maketitle

\else
\fi

% Theorems
\theoremstyle{definition}
\newtheorem*{aim}{Aim}
\newtheorem{axiom}{Axiom}[subsection]
\newtheorem*{claim}{Claim}
\newtheorem*{corollary}{Corollary}
\newtheorem*{conjecture}{Conjecture}
\newtheorem{definition}{Definition}[subsection]
\newtheorem*{eg}{Example}
\newtheorem*{ex}{Exercise}
\newtheorem*{fact}{Fact}
\newtheorem*{law}{Law}
\newtheorem{lemma}{Lemma}[subsubsection]
\newtheorem*{notation}{Notation}
\newtheorem{proposition}{Proposition}[definition]
\newtheorem{property}{Property}[subsection]
\newtheorem*{question}{Question}
\newtheorem*{problem}{Problem}
\newtheorem*{rrule}{Rule}
\newtheorem{theorem}{Theorem}[subsection]
\newtheorem*{assumption}{Assumption}

\newtheorem*{remark}{Remark}
\newtheorem*{warning}{Warning}
\newtheorem*{exercise}{Exercise}

\newtheorem{nthm}{Theorem}[section]
\newtheorem{nlemma}[nthm]{Lemma}
\newtheorem{nprop}[nthm]{Proposition}
\newtheorem{ncor}[nthm]{Corollary}


\renewcommand{\labelitemi}{--}
\renewcommand{\labelitemii}{$\circ$}
\renewcommand{\labelenumi}{(\roman{*})}

\let\stdsection\section
\renewcommand\section{\newpage\stdsection}




%%%%%%%%%%%%%%%%%%%%%%%%%
%%%%% Maths Symbols %%%%%
%%%%%%%%%%%%%%%%%%%%%%%%%

\let\U\relax
\let\C\relax
\let\G\relax

% Matrix groups
\newcommand{\GL}{\mathrm{GL}}
\newcommand{\Or}{\mathrm{O}}
\newcommand{\PGL}{\mathrm{PGL}}
\newcommand{\PSL}{\mathrm{PSL}}
\newcommand{\PSO}{\mathrm{PSO}}
\newcommand{\PSU}{\mathrm{PSU}}
\newcommand{\SL}{\mathrm{SL}}
\newcommand{\SO}{\mathrm{SO}}
\newcommand{\Spin}{\mathrm{Spin}}
\newcommand{\Sp}{\mathrm{Sp}}
\newcommand{\SU}{\mathrm{SU}}
\newcommand{\U}{\mathrm{U}}
\newcommand{\Mat}{\mathrm{Mat}}

% Matrix algebras
\newcommand{\gl}{\mathfrak{gl}}
\newcommand{\ort}{\mathfrak{o}}
\newcommand{\so}{\mathfrak{so}}
\newcommand{\su}{\mathfrak{su}}
\newcommand{\uu}{\mathfrak{u}}
\renewcommand{\sl}{\mathfrak{sl}}

% Special sets
\newcommand{\C}{\mathbb{C}}
\newcommand{\CP}{\mathbb{CP}}
\newcommand{\GG}{\mathbb{G}}
\newcommand{\N}{\mathbb{N}}
\newcommand{\Q}{\mathbb{Q}}
\newcommand{\R}{\mathbb{R}}
\newcommand{\RP}{\mathbb{RP}}
\newcommand{\T}{\mathbb{T}}
\newcommand{\Z}{\mathbb{Z}}
\renewcommand{\H}{\mathbb{H}}

% Brackets
\newcommand{\abs}[1]{\left\lvert #1\right\rvert}
\newcommand{\bket}[1]{\left\lvert #1\right\rangle}
\newcommand{\brak}[1]{\left\langle #1 \right\rvert}
\newcommand{\braket}[2]{\left\langle #1\middle\vert #2 \right\rangle}
\newcommand{\bra}{\langle}
\newcommand{\ket}{\rangle}
\newcommand{\norm}[1]{\left\lVert #1\right\rVert}
\newcommand{\normalorder}[1]{\mathop{:}\nolimits\!#1\!\mathop{:}\nolimits}
\newcommand{\tv}[1]{|#1|}
\renewcommand{\vec}[1]{\boldsymbol{\mathbf{#1}}}

% not-math
\newcommand{\bolds}[1]{{\bfseries #1}}
\newcommand{\cat}[1]{\mathsf{#1}}
\newcommand{\ph}{\,\cdot\,}
\newcommand{\term}[1]{\emph{#1}\index{#1}}
\newcommand{\phantomeq}{\hphantom{{}={}}}
% Probability
\DeclareMathOperator{\Bernoulli}{Bernoulli}
\DeclareMathOperator{\betaD}{beta}
\DeclareMathOperator{\bias}{bias}
\DeclareMathOperator{\binomial}{binomial}
\DeclareMathOperator{\corr}{corr}
\DeclareMathOperator{\cov}{cov}
\DeclareMathOperator{\gammaD}{gamma}
\DeclareMathOperator{\mse}{mse}
\DeclareMathOperator{\multinomial}{multinomial}
\DeclareMathOperator{\Poisson}{Poisson}
\DeclareMathOperator{\var}{var}
\newcommand{\E}{\mathbb{E}}
\newcommand{\Prob}{\mathbb{P}}

% Algebra
\DeclareMathOperator{\adj}{adj}
\DeclareMathOperator{\Ann}{Ann}
\DeclareMathOperator{\Aut}{Aut}
\DeclareMathOperator{\Char}{char}
\DeclareMathOperator{\disc}{disc}
\DeclareMathOperator{\dom}{dom}
\DeclareMathOperator{\fix}{fix}
\DeclareMathOperator{\Hom}{Hom}
\DeclareMathOperator{\id}{id}
\DeclareMathOperator{\image}{image}
\DeclareMathOperator{\im}{im}
\DeclareMathOperator{\re}{re}
\DeclareMathOperator{\tr}{tr}
\DeclareMathOperator{\Tr}{Tr}
\newcommand{\Bilin}{\mathrm{Bilin}}
\newcommand{\Frob}{\mathrm{Frob}}

% Others
\newcommand\ad{\mathrm{ad}}
\newcommand\Art{\mathrm{Art}}
\newcommand{\B}{\mathcal{B}}
\newcommand{\cU}{\mathcal{U}}
\newcommand{\Der}{\mathrm{Der}}
\newcommand{\D}{\mathrm{D}}
\newcommand{\dR}{\mathrm{dR}}
\newcommand{\exterior}{\mathchoice{{\textstyle\bigwedge}}{{\bigwedge}}{{\textstyle\wedge}}{{\scriptstyle\wedge}}}
\newcommand{\F}{\mathbb{F}}
\newcommand{\G}{\mathcal{G}}
\newcommand{\Gr}{\mathrm{Gr}}
\newcommand{\haut}{\mathrm{ht}}
\newcommand{\Hol}{\mathrm{Hol}}
\newcommand{\hol}{\mathfrak{hol}}
\newcommand{\Id}{\mathrm{Id}}
\newcommand{\lie}[1]{\mathfrak{#1}}
\newcommand{\op}{\mathrm{op}}
\newcommand{\Oc}{\mathcal{O}}
\newcommand{\pr}{\mathrm{pr}}
\newcommand{\Ps}{\mathcal{P}}
\newcommand{\pt}{\mathrm{pt}}
\newcommand{\qeq}{\mathrel{``{=}"}}
\newcommand{\Rs}{\mathcal{R}}
\newcommand{\Vect}{\mathrm{Vect}}
\newcommand{\wsto}{\stackrel{\mathrm{w}^*}{\to}}
\newcommand{\wt}{\mathrm{wt}}
\newcommand{\wto}{\stackrel{\mathrm{w}}{\to}}
\renewcommand{\d}{\mathrm{d}}
\renewcommand{\P}{\mathbb{P}}
%\renewcommand{\F}{\mathcal{F}}


\let\Im\relax
\let\Re\relax

\DeclareMathOperator{\area}{area}
\DeclareMathOperator{\card}{card}
\DeclareMathOperator{\ccl}{ccl}
\DeclareMathOperator{\ch}{ch}
\DeclareMathOperator{\cl}{cl}
\DeclareMathOperator{\cls}{\overline{\mathrm{span}}}
\DeclareMathOperator{\coker}{coker}
\DeclareMathOperator{\conv}{conv}
\DeclareMathOperator{\cosec}{cosec}
\DeclareMathOperator{\cosech}{cosech}
\DeclareMathOperator{\covol}{covol}
\DeclareMathOperator{\diag}{diag}
\DeclareMathOperator{\diam}{diam}
\DeclareMathOperator{\Diff}{Diff}
\DeclareMathOperator{\End}{End}
\DeclareMathOperator{\energy}{energy}
\DeclareMathOperator{\erfc}{erfc}
\DeclareMathOperator{\erf}{erf}
\DeclareMathOperator*{\esssup}{ess\,sup}
\DeclareMathOperator{\ev}{ev}
\DeclareMathOperator{\Ext}{Ext}
\DeclareMathOperator{\fst}{fst}
\DeclareMathOperator{\Fit}{Fit}
\DeclareMathOperator{\Frac}{Frac}
\DeclareMathOperator{\Gal}{Gal}
\DeclareMathOperator{\gr}{gr}
\DeclareMathOperator{\hcf}{hcf}
\DeclareMathOperator{\Im}{Im}
\DeclareMathOperator{\Ind}{Ind}
\DeclareMathOperator{\Int}{Int}
\DeclareMathOperator{\Isom}{Isom}
\DeclareMathOperator{\lcm}{lcm}
\DeclareMathOperator{\length}{length}
\DeclareMathOperator{\Lie}{Lie}
\DeclareMathOperator{\like}{like}
\DeclareMathOperator{\Lk}{Lk}
\DeclareMathOperator{\Maps}{Maps}
\DeclareMathOperator{\orb}{orb}
\DeclareMathOperator{\ord}{ord}
\DeclareMathOperator{\otp}{otp}
\DeclareMathOperator{\poly}{poly}
\DeclareMathOperator{\rank}{rank}
\DeclareMathOperator{\rel}{rel}
\DeclareMathOperator{\Rad}{Rad}
\DeclareMathOperator{\Re}{Re}
\DeclareMathOperator*{\res}{res}
\DeclareMathOperator{\Res}{Res}
\DeclareMathOperator{\Ric}{Ric}
\DeclareMathOperator{\rk}{rk}
\DeclareMathOperator{\Rees}{Rees}
\DeclareMathOperator{\Root}{Root}
\DeclareMathOperator{\sech}{sech}
\DeclareMathOperator{\sgn}{sgn}
\DeclareMathOperator{\snd}{snd}
\DeclareMathOperator{\Spec}{Spec}
\DeclareMathOperator{\spn}{span}
\DeclareMathOperator{\stab}{stab}
\DeclareMathOperator{\St}{St}
\DeclareMathOperator{\supp}{supp}
\DeclareMathOperator{\Syl}{Syl}
\DeclareMathOperator{\Sym}{Sym}
\DeclareMathOperator{\vol}{vol}

\pgfarrowsdeclarecombine{twolatex'}{twolatex'}{latex'}{latex'}{latex'}{latex'}
\tikzset{->/.style = {decoration={markings,
                                  mark=at position 1 with {\arrow[scale=2]{latex'}}},
                      postaction={decorate}}}
\tikzset{<-/.style = {decoration={markings,
                                  mark=at position 0 with {\arrowreversed[scale=2]{latex'}}},
                      postaction={decorate}}}
\tikzset{<->/.style = {decoration={markings,
                                   mark=at position 0 with {\arrowreversed[scale=2]{latex'}},
                                   mark=at position 1 with {\arrow[scale=2]{latex'}}},
                       postaction={decorate}}}
\tikzset{->-/.style = {decoration={markings,
                                   mark=at position #1 with {\arrow[scale=2]{latex'}}},
                       postaction={decorate}}}
\tikzset{-<-/.style = {decoration={markings,
                                   mark=at position #1 with {\arrowreversed[scale=2]{latex'}}},
                       postaction={decorate}}}
\tikzset{->>/.style = {decoration={markings,
                                  mark=at position 1 with {\arrow[scale=2]{latex'}}},
                      postaction={decorate}}}
\tikzset{<<-/.style = {decoration={markings,
                                  mark=at position 0 with {\arrowreversed[scale=2]{twolatex'}}},
                      postaction={decorate}}}
\tikzset{<<->>/.style = {decoration={markings,
                                   mark=at position 0 with {\arrowreversed[scale=2]{twolatex'}},
                                   mark=at position 1 with {\arrow[scale=2]{twolatex'}}},
                       postaction={decorate}}}
\tikzset{->>-/.style = {decoration={markings,
                                   mark=at position #1 with {\arrow[scale=2]{twolatex'}}},
                       postaction={decorate}}}
\tikzset{-<<-/.style = {decoration={markings,
                                   mark=at position #1 with {\arrowreversed[scale=2]{twolatex'}}},
                       postaction={decorate}}}

\tikzset{circ/.style = {fill, circle, inner sep = 0, minimum size = 3}}
\tikzset{scirc/.style = {fill, circle, inner sep = 0, minimum size = 1.5}}
\tikzset{mstate/.style={circle, draw, blue, text=black, minimum width=0.7cm}}

\tikzset{eqpic/.style={baseline={([yshift=-.5ex]current bounding box.center)}}}
\tikzset{commutative diagrams/.cd,cdmap/.style={/tikz/column 1/.append style={anchor=base east},/tikz/column 2/.append style={anchor=base west},row sep=tiny}}

\definecolor{mblue}{rgb}{0.2, 0.3, 0.8}
\definecolor{morange}{rgb}{1, 0.5, 0}
\definecolor{mgreen}{rgb}{0.1, 0.4, 0.2}
\definecolor{mred}{rgb}{0.5, 0, 0}

\def\drawcirculararc(#1,#2)(#3,#4)(#5,#6){%
    \pgfmathsetmacro\cA{(#1*#1+#2*#2-#3*#3-#4*#4)/2}%
    \pgfmathsetmacro\cB{(#1*#1+#2*#2-#5*#5-#6*#6)/2}%
    \pgfmathsetmacro\cy{(\cB*(#1-#3)-\cA*(#1-#5))/%
                        ((#2-#6)*(#1-#3)-(#2-#4)*(#1-#5))}%
    \pgfmathsetmacro\cx{(\cA-\cy*(#2-#4))/(#1-#3)}%
    \pgfmathsetmacro\cr{sqrt((#1-\cx)*(#1-\cx)+(#2-\cy)*(#2-\cy))}%
    \pgfmathsetmacro\cA{atan2(#2-\cy,#1-\cx)}%
    \pgfmathsetmacro\cB{atan2(#6-\cy,#5-\cx)}%
    \pgfmathparse{\cB<\cA}%
    \ifnum\pgfmathresult=1
        \pgfmathsetmacro\cB{\cB+360}%
    \fi
    \draw (#1,#2) arc (\cA:\cB:\cr);%
}
\newcommand\getCoord[3]{\newdimen{#1}\newdimen{#2}\pgfextractx{#1}{\pgfpointanchor{#3}{center}}\pgfextracty{#2}{\pgfpointanchor{#3}{center}}}

\newcommand\qedshift{\vspace{-17pt}}
\newcommand\fakeqed{\pushQED{\qed}\qedhere}

\def\Xint#1{\mathchoice
   {\XXint\displaystyle\textstyle{#1}}%
   {\XXint\textstyle\scriptstyle{#1}}%
   {\XXint\scriptstyle\scriptscriptstyle{#1}}%
   {\XXint\scriptscriptstyle\scriptscriptstyle{#1}}%
   \!\int}
\def\XXint#1#2#3{{\setbox0=\hbox{$#1{#2#3}{\int}$}
     \vcenter{\hbox{$#2#3$}}\kern-.5\wd0}}
\def\ddashint{\Xint=}
\def\dashint{\Xint-}

\newcommand\separator{{\centering\rule{2cm}{0.2pt}\vspace{2pt}\par}}

\newenvironment{own}{\color{gray!70!black}}{}

\newcommand\makecenter[1]{\raisebox{-0.5\height}{#1}}

\mathchardef\mdash="2D

\newenvironment{significant}{\begin{center}\begin{minipage}{0.9\textwidth}\centering\em}{\end{minipage}\end{center}}
\DeclareRobustCommand{\rvdots}{%
  \vbox{
    \baselineskip4\p@\lineskiplimit\z@
    \kern-\p@
    \hbox{.}\hbox{.}\hbox{.}
  }}
\DeclareRobustCommand\tph[3]{{\texorpdfstring{#1}{#2}}}

\makeatother

\newcommand{\nn}{(natural\;number)}
\newcommand{\at}{\text{Assum that }}
\newcommand{\exforall}{\forall\;}
\newcommand{\st}{,\;\operatorname{s.t.}\;}
\newcommand{\exexists}{\exists\;}
\begin{document}
\everymath{\displaystyle}
\maketitle
\tableofcontents
\setcounter{chapter}{-1}
\chapter{The construction of the real number}
\paragraph{}
The chapter mainly defines the real and the set, which is the basis of Analysis.
\paragraph{}
The chapter follows the natural numbers-rational number-real numbers routine to construct the real number.
\section{Natural number}
\subsection{The construction of the natural number}
\paragraph{Peano axioms}
\begin{definition}
    $a++$ means the subsequence of $a$.
\end{definition}
\begin{definition}
    $(x)m$ means $m$ is belong to $x$ or $m$ is $x$.
    
    I.e.,$\;\nn x$ means $x$ is a natural number. 
\end{definition}
\begin{definition}
    "$a\coloneqq b$" means $a$ is defined by $b$,means $a\Leftrightarrow b$ 
\end{definition}
\begin{axiom}
    $$\nn 0\label{axiom1111}$$
\end{axiom}
\begin{axiom}
    $$\nn n\Rightarrow\nn (n++)\label{axiom1112}$$
\end{axiom}
\begin{definition}
    $1\coloneqq 0++\quad,\;2\coloneqq (0++)++\quad,\;\cdots$
\end{definition}
\begin{axiom}
    $$\exforall\nn n,(n++)\neq0\label{axiom1113}$$
\end{axiom}
\begin{axiom}
    $$ [\nn b\neq \nn c\Rightarrow(b++)\neq(c++)]$$
    $$\Leftrightarrow[(b++)=(c++)\Rightarrow(b=c)]\label{axiom1114}$$
\end{axiom}
\begin{axiom}[Principle of mathematical induction]
    $$\{P\nn m_0),[P((natural\;number,\;>m_0)n)$$
    $$\Rightarrow P(n++)]\}\Rightarrow\exforall\nn n,P(n),n\geq m_0\label{axiom1115}$$
\end{axiom}
\subsection{The operation of natural numbers}
\subsubsection{Addition}
\begin{assumption}
    $ $

    $\text{(unofficial)} \exexists \text{a number system }\mathbb{N},\;\text{the element of }\mathbb{N}\text{ is called natural}$

    $ \text{number,and Axiom }\ref{axiom1111}\;to\;Axiom\;\ref{axiom1115}\text{ holds for }\mathbb{N}.$

    tip:The assumption can be more precise after the set being defined.
\end{assumption}
\begin{proposition}[recursive definition]
    $$\exforall \nn n,\exexists f_n:\,\mathbb{N}\mapsto\mathbb{N},\;a_0=c,\;a_{n+1}=f_n(a_n)$$
\end{proposition}
\begin{proof}
    $$(informal)a_0=c,\;\exforall\nn n,\;c\text{ isn't defined by }f_n\;$$$$\text{according to }Axiom\;\ref{axiom1113}.$$
    $$\at a_n=m,a_{n++}\coloneqq f_n(a_n)\text{, according to }Axiom\;\ref{axiom1115},$$
    $$\exforall \nn n,a_{n++}\coloneqq f_n(a_n)$$
\end{proof}
\begin{definition}
    $$0+m\coloneqq m$$
\end{definition}
\begin{definition}
    $$(n++)+m\coloneqq(n+m)++$$
\end{definition}
\begin{proposition}
    $$\nn n,\;\nn m\Rightarrow\nn (n+m)$$
\end{proposition}
\begin{proof}
    $$0+\nn n=n$$
    $$\at \nn m+n=\nn t$$
    $$\therefore (m++)+n=(m+n)++=t++$$
    $$\because \nn t,\;\therefore\nn(t++)$$

    $$(\nn a++)+\nn b=(a+b)++$$
\end{proof}
\begin{lemma}
    $$\exforall\nn n,n+0=n$$
\end{lemma}
\begin{proof}
    $$0+0=0$$
    $$\at n+0=n$$
    $$(n++)+0=(n+0)++=n++$$
\end{proof}
\begin{lemma}
    $$\exforall\nn n,\exforall\nn m,n+(m++)=(n+m)++$$
\end{lemma}
\begin{proof}
    $$0+(\nn m++)=(0+m)++$$
    $$\at \nn n+(m++)=(n+m)++$$
    $$(n+m)++=(n++)+m$$
    $$\therefore\;(n++)+(m++)=(n+(m++))++=((n+m)++)++$$
    $$\exforall\nn n,n+(\nn m++)=(n+m)++$$
    $$$$
    $$\nn n+0=0$$
    $$\at n+(\nn m++)=(n+m)++$$
    $$(n+m)++=(n++)+m$$
    $$\therefore\;(n++)+(m++)=((n++)+m)++=(n+(m++))++=n+((m++)++)$$
    $$\exforall\nn m,\nn n+(m++)=(n+m)++$$
\end{proof}
\begin{corollary}
    $\nn n++=n+1$
\end{corollary}
\begin{proof}
    $\nn n++=(n+0)++=n+(0++)=n+1$
\end{proof}
\begin{theorem}[Addition is commutative]
    $$\exforall\nn n,\exforall\nn m,n+m=m+n$$
\end{theorem}
\begin{proof}
    $$\nn n+0=n=0+n$$
    $$\at n+\nn m=m+n$$
    $$(n++)+m=(n+m)++=(n+m)++=(m+n)++=m+(n++)$$
    $$\exforall\nn m,\;\nn n+m=m+n$$
    $$$$
    $$0+\nn m=m=m+0$$
    $$\at \nn n+m=m+n$$
    $$n+(m++)=(n+m)++=(m+n)++=(m++)+n$$
    $$\exforall\nn n,\;n+\nn m=m+n$$
\end{proof}
\begin{theorem}[Addition is associative]
    $$(\nn a+\nn b)+\nn c=a+(b+c)$$
\end{theorem}
\begin{proof}
    $$\begin{aligned}
        (\nn a+\nn b)+\nn c=a+b+c\\
        =b+a+c=b+c+a=(b+c)+a
    \end{aligned}$$
\end{proof}
\begin{theorem}[Cancellation law]
    $$\nn a+\nn b=a+\nn c\Rightarrow b=c$$
\end{theorem}
\begin{proof}
    $$0+\nn b=0+\nn c\Rightarrow b=c$$
    $$\at \nn a+b=a+c\Rightarrow b=c$$
    $$(a++)+b=(a+b)++,\;(a++)+c=(a+c)++$$
    $$\because a+b=a+c$$
    $$\therefore (a+b)++=(a+c)++$$
\end{proof}
\subsubsection{Natural number sequence}
\begin{definition}
    $\nn n\text{ is positive}\Leftrightarrow n\neq 0$
\end{definition}
\begin{theorem}
    $(positive)a,\;\nn b,\;(positive)(a+b)$
\end{theorem}
\begin{proof}
    $$a+0=a$$
    $$\at (positive)(a+\nn b)$$
    $$a+(b++)=(a+b)++$$
    $$\therefore (positive)(a+(b++))$$
\end{proof}
\begin{theorem}
    $\nn a+\nn b=0\Leftrightarrow a=b=0$
\end{theorem}
\begin{proof}
    $$\at\;a\neq0,b\neq0,\st a+b=0$$
    $$\text{Obviously},\exforall(positive)a,\exforall(positive)b,\;a+b\neq0$$
\end{proof}
\begin{lemma}
    $\exforall(positive)n,\exexists\nn a\st a++=n$
\end{lemma}
\begin{proof}
    $$n=1,a_1=0$$
    $$\at n=k,\;\exexists a_n++=k$$
    $$when\;n=k++,a_{n+1}=a_{n}++=k$$
\end{proof}
\begin{definition}[Ordering of the natural numbers]
    $$\nn n\geq\nn m\Leftrightarrow m\leq n$$
    $$\Leftrightarrow\exexists\nn a\st n=m+a$$
    $$n>m\Leftrightarrow m<n\Leftrightarrow\exexists(positive) a\st n=m+a$$
\end{definition}
\begin{proposition}[Basic properties of order for natural numbers]
    $$\nn a,\;\nn b,\;\nn c$$
    $$\begin{aligned}
        &a)\qquad(Order\;is\;reflexive)\qquad a\geq a\\
        &b)\qquad(Order\;is\;transitive)\qquad (a\geq b,\;b\geq c)\Rightarrow a\geq c\\
        &c)\qquad(Order\;is\;anti-symmetric)\qquad (a\geq b,\;a\leq b)\Rightarrow a=b\\
        &d)\qquad(Addition\;preserves\;order)\qquad a\geq b\Leftrightarrow a+c\geq b+c\\
        &e)\qquad a<b\Leftrightarrow a++\geq b\\
        &f)\qquad a<b\Leftrightarrow \exexists(positive)d\st b=a+d
    \end{aligned}$$
\end{proposition}
\begin{proof}
    $$\begin{aligned}
        &a)\quad a=a+0\Rightarrow a\geq a\\
        &b)\quad (a=b+\nn x,b=c+\nn y)\\
        &\qquad \Rightarrow a=c+x+y\geq c\\
        &c)\quad (a\geq b\Rightarrow a=b+\nn x,a\leq b\Rightarrow a+\nn y=b)\\
        &\qquad\therefore a=a+x+y\Rightarrow x+y=0\Rightarrow x=y=0\\
        &d)\quad a\geq b\Rightarrow a=b+\nn x\\
        &\qquad a+0=b+x+0\Rightarrow a+0\geq b+0\\
        &\qquad \at a+\nn c=b+c+x\Rightarrow a+c\geq b+c\\
        &\qquad \therefore a+(c++)=(a+c)++=((b+c)+x)++=b+(c++)+x\\
        &\qquad \Rightarrow a+c\geq b+c\\
        &e)\quad a+(positive)x=b,\;obviously\;x\geq1\Rightarrow a+1\leq b\\
        &f)\quad b=a+1\Rightarrow d=1\\
        &\qquad\at b=a+(positive)n\Rightarrow d=n\\
        &\qquad b++=a+(n++)\Rightarrow d=n++\\
    \end{aligned}$$
\end{proof}
\begin{proposition}[Trichotomy of order for natural numbers]
    $$\text{There'll there exists unique one statement is true : }a<b,\;a=b,\;a>b$$
\end{proposition}
\begin{proof}
    $$\text{First, prove no more than one of the statements holding at the same time.}$$
    $$a<b\Rightarrow a+(positive)x=b\Rightarrow a\neq b,\;a\not\geq b$$
    $$a=b\Rightarrow \exforall(positive)x,\;x+a>b,\;a<b+x\Rightarrow a\not>b,a\not<b$$
    $$a>b\Rightarrow a=b+(positive)x\Rightarrow a\neq b,\;a\not\leq b$$

    $$\text{Then prove no less than one of the statements holding at the same time.}$$
    $$P(\nn x)\Leftrightarrow\{\exforall x,\exforall \nn y$$
    $$\rightarrow[(x<y)\vee(x=y)\vee(x<y)]\}$$
    $$P(0)\Leftrightarrow 0\leq y\Leftrightarrow(0<y)\vee(0=y)$$
    $$Obivously,\;\not\exexists(positive)z,\;0+z=y\st 0=0+z$$
    $$\at P(x)$$
    $$if\;x<y\Rightarrow (x++)\leq y\Rightarrow((x++)<y)\vee((x++)=y)$$
    $$Obivously,\;\not\exexists(positive)z,\;x+z=y\st x=x+z$$
    $$if\;x\geq y\Rightarrow (x++)>y$$
    $$\therefore \exforall\nn x,P(x)$$
\end{proof}
\begin{theorem}[Strong principle of induction]
    $$\{P(\nn m_0),[[P(0)\vee\cdots\vee P((natural\;number,\;\geq m_0)n)]$$
    $$\Rightarrow P(n++)\}\Rightarrow P(n)$$
\end{theorem}
\begin{proof}
    $$Construct\;G(n)=P(m_0)\vee\cdots\vee P(n)$$
    $$P(0)\Rightarrow G(0)$$
    $$\at G(n)\Rightarrow P(n)$$
    $$[G(n)\Rightarrow G(n++)]\Rightarrow P(n++))$$
\end{proof}
\begin{theorem}[principle of backwards induction]
    $$\{P(\nn n),[P(((natural\;number,\;\leq n)m)++)\Rightarrow P(m)]\}$$
    $$\Rightarrow P(m)$$
\end{theorem}
\begin{proof}
    $$Construct\;G(\nn a)=P(0)\vee\cdots\vee P(a)$$
    $$Obviously,\;P(0)\Rightarrow Q(0)$$
    $$\at Q(n)=P(0)\vee\cdots\vee P(n)$$
    $$P(0)\vee\cdots\vee P(n)\Rightarrow\text{[Strong principle of induction]}P(n++)\Rightarrow G(n++)$$
\end{proof}
\subsubsection{Multiplication}
\begin{definition}
    $0\times m\coloneqq0$
\end{definition}
\begin{definition}
    $(n++)\times m\coloneqq(n\times m)+m$
\end{definition}
\begin{lemma}
    $m\times 0=0$
\end{lemma}
\begin{proof}
    $$0\times0=0$$
    $$\at n\times0=0$$
    $$(n++)\times0=(n\times 0)+0=0$$
\end{proof}
\begin{lemma}
    $\nn n\times(\nn m++)=(n\times m)+n$
\end{lemma}
\begin{proof}
    $$0\times(\nn m++)=(m++)\times 0=0$$
    $$\at (\nn n)\times (m++)=(n\times m)+n$$
    $$(n++)\times(m++)=(n\times(m++))+(m++)=n\times m+(m++)+n$$
    $$\therefore\;(n++)\times(m++)=n\times m+m+(n++)=(n\times(m++))+(n++)$$
\end{proof}
\begin{theorem}[Multiplication is commutative]
    $$m\times n=n\times m$$
\end{theorem}
\begin{proof}
    $$0\times (\nn m++)=(m++)\times0=0$$
    $$\at \nn n\times m=m\times n$$
    $$(n++)\times m=n\times m+m=m\times n+m=(m++)\times n$$
\end{proof}
\end{document}